\chapter*{Introduction Générale }
\label{Intro}
\markboth{Introduction Générale}{Introduction Générale}
\mtcaddchapter[Introduction Générale]

\minitoc


À l’ère du numérique et de la transformation digitale, les technologies mobiles, l’Internet des objets et l’intelligence artificielle occupent une place centrale dans l’amélioration de la qualité de vie et de la sécurité des individus. Les smartphones, les systèmes de géolocalisation en temps réel et les plateformes connectées permettent aujourd’hui de surveiller, d’analyser et de réagir rapidement à des situations critiques. Dans ce contexte, la sécurité des personnes vulnérables, notamment les enfants, les personnes âgées et celles atteintes de troubles cognitifs, représente un enjeu sociétal majeur.

En effet, ces catégories de personnes sont particulièrement exposées à des risques tels que la perte d’orientation, les fugues, les accidents ou les situations d’urgence médicale. Les familles et les organisations responsables de leur prise en charge font face à des difficultés importantes pour assurer un suivi continu, réagir rapidement en cas de danger et disposer d’informations fiables sur leurs déplacements. Bien que plusieurs solutions existent sur le marché, elles restent souvent limitées, peu personnalisables ou ne proposent pas une approche globale intégrant à la fois le suivi en temps réel, la gestion des alertes et l’analyse intelligente des comportements.

C’est dans ce cadre que s’inscrit le projet SafeTrack. SafeTrack est une application mobile de suivi et de sécurité destinée aux personnes vulnérables, offrant une solution complète et intelligente pour la surveillance, la prévention des risques et la gestion des situations d’urgence. Basée sur la géolocalisation en temps réel, l’historique des déplacements et la définition de zones sécurisées (géofencing), l’application permet d’alerter instantanément les responsables en cas de comportement anormal ou de dépassement de zones autorisées. Elle intègre également un bouton SOS et des fonctionnalités d’appels d’urgence afin d’assurer une intervention rapide en situation critique.

SafeTrack se distingue par l’intégration d’un dashboard web sécurisé destiné aux organisations et aux responsables, ainsi que par l’exploitation de l’intelligence artificielle, notamment à travers des mécanismes de détection intelligente des comportements à risque. Cette approche permet non seulement un suivi passif, mais également une anticipation des situations dangereuses, contribuant ainsi à renforcer la prévention et la prise de décision.
% non precis (selon le nbr des chapitres)
Ce rapport est structuré en plusieurs chapitres. Le premier chapitre présente le contexte général du projet, la problématique abordée, ainsi qu’une analyse des solutions existantes et de la valeur ajoutée de SafeTrack. Le deuxième chapitre est consacré à l’étude des besoins fonctionnels et non fonctionnels, à la modélisation du système et à la présentation de l’architecture globale et des technologies utilisées.

% Le deuxième chapitre est consacré aux besoins fonctionnels et non fonctionnels du projet. Il inclut le backlog général, les diagrammes de cas d’utilisation ainsi que le diagramme de classes général, accompagnés d’une explication de l’architecture de notre plateforme. Ce chapitre aborde aussi les outils et l’environnement de travail utilisés.

% Le troisième chapitre décrit la réalisation des deux premiers sprints : le sprint un, portant sur l’authentification et la gestion des utilisateurs pour l’administrateur, et le sprint deux, centré sur la gestion des profils et les relations de suivi entre utilisateurs.

% Le quatrième chapitre couvre les sprints trois et quatre. Le sprint trois traite de la gestion des posts, des stories et des interactions sociales, tandis que le sprint quatre est dédié au tableau de bord administrateur ainsi qu’au système de notifications.

% Le cinquième chapitre détaille le sprint cinq, qui inclut la gestion des offres d’emploi, des projets freelance, ainsi que le système de messagerie.

% Le sixième chapitre est consacré au sprint six, qui introduit le système de réunions vidéo.

% Enfin, le septième chapitre présente le sprint sept, intégrant les fonctionnalités d’intelligence artificielle.

% Le rapport se conclut par une synthèse des résultats obtenus, une évaluation de l’impact du projet, et des recommandations pour les améliorations et évolutions futures




